\section{Conclusion}
\label{sec:conclusion}

Solving a small differentiable constrained projection at every decoding step, and training
through it, produces a model whose samples individually respect a quantitative target far
more often than samples from the distribution the field treats as optimal for that target.
That inversion is the paper's substantive finding, and it has a mechanism rather than a
coincidence behind it: an expectation constraint is indifferent to individual draws, so the
cheapest route to a mean is through the tails, while a per-step feasible set is centred on
what the remaining budget can absorb. Because the target enters as a value rather than a
training condition, the same trained network serves an entire family of targets, which
relocates the cost of control from inference to a single training run.

\paragraph{Limitations}
The projection reweights the backbone's own proposals, so it cannot reach a target the
backbone assigns no mass to, and the achievable range saturates well before the requested
one does. Two correlated constraints sharing a single token budget compose only partially,
and asymmetrically: the direction that spends less of the budget is reliably available
while the direction that spends more is not. The expressible class is fixed by the
serialization, since a target no prefix determines cannot be steered at all, and charge
neutrality falls outside that class rather than merely being unimplemented. On accuracy per
target the method is not the strongest option available, because sampling repeatedly from
an unsteered model and selecting the closest output wins given enough inference compute.
The synthetic results are exact but small, and the crystal energetics rest on a surrogate
rather than density functional theory.

\paragraph{Future work}
Three threads follow directly. The two-constraint grid is one row short of the range needed
to separate a smooth dependence on target compatibility from a jump, and completing it is a
matter of compute rather than method. The particle-filter comparison used an untuned
potential width, so the figure reported here bounds that family from below rather than
characterising it, and a tuned version would sharpen the comparison the paper draws. The
serialization result is currently a statement about two orderings measured on one quantity;
turning it into a predictive account of which constraints a given serialization admits needs
several serializations budgeted together, which is the form in which it would become a
design rule rather than an observation.

%%%%%%%%%%%%%%%%%%%%%%%%%%%%%%%%%%%%%%%%%%%%%%%%%%%%%%%%%%%%%%%%%%%%%%%%%%%%%%%%
